%!TEX root = ../main.tex

\section{Related Work}
\paragraph{Pruning.}
Structured pruning has been extensively studied as a model compression technique in computer vision and natural language processing, 
where task-specific models like  classification ones are often overparameterized and can be pruned significantly with minimal impact on performance~\citep{han2016deep, wen2016learning,liu2017learning,luo2017thinet, cai2019once, deng2020model,hou2020dynabert, wang2020structured, lagunas2021block, xia-etal-2022-structured, kurtic2023ziplm}. 
Unstructured pruning~\citep{frankle2018lottery,li2020train, chen2020lottery, sanh2020movement} prunes individual neurons instead of structured blocks.
Though unstructured pruning usually achieve higher compression rates,
they are  not practical for model speedup.

In the era of LLMs, 
the prevalent NLP pipeline has shifted from task-specific models 
to general-purpose LMs,
which leaves little room for redundancy. 
Both unstructured pruning, 
semi-structured pruning~\citep{frantar2023sparsegpt,sun2023simple}, 
and structured pruning~\citep{ma2023llm}
lead to significant performance drops on LLM even at a modest sparsity.
Noticeably, all previous works 
fix the original models or tune them minimally.
We
see pruning as an initialization and 
consider it necessary to expend substantial compute to 
continually pre-training the model to recover performance. 

\paragraph{Efficient pre-training approaches.}
As orthogonal to our pruning approach, 
There is an extensive body of work on improving efficiency of training LLMs. 
For example, 
quantization reduces the numeric precision of  model weights and activations 
and speeds up training and inference~\citep{dettmers2022llm,dettmers2023qlora,xiao2023smoothquant}.
Knowledge distillation~\citep{hinton2015distilling,sanh2019distilbert,jiao2020tinybert,sun2020mobilebert}, 
which trains a smaller model on a larger model's prediction,
is shown to be effective for task-specific models~\citep{xia-etal-2022-structured}. For pre-training LLMs, though distilling from a teacher model is shown to improve the quality of student models given the same number of training steps~\citep{rae2021scaling, blakeney2022reduce}, it is less cost-effective than pruning and continued training due to the exceeding inference cost incured by the teacher model~\citep{jha2023large}. More  methods have been introduced to enhance the efficiency of training LMs, 
such as dynamic architectures~\citep{gong2019efficient,zhang2020accelerating} and efficient optimizers~\citep{chen2023symbolic, liu2023sophia}. However, 
as indicated by~\citep{kaddour2023no,bartoldson2023compute},
the promised gains in training efficiency may not be consistently realized.

There are also data-based approaches to enhance training efficiency.
Eliminating duplicated data is found to be effective
\citep{lee2021deduplicating}. 
Various batch selection techniques propose to prioritize data based on criteria such as 
higher losses~\citep{jiang2019accelerating} or a greater reducible loss~\citep{mindermann2022prioritized}. 
\cite{xie2023doremi}
propose to 
optimize data mixtures 
by training a proxy model to estimate the optimal data weight of each domain. 






